\documentclass{szb-solution}

\area{Integrálszámítás}
\title{Integrálszámítás II}
\subject{Matematika G1}
\subjectCode{BMETE94BG01}
\date{Utoljára frissítve: \today}
\docno{11}

\begin{document}
\maketitle

\subsection{Parciális integrálás}

Minden eredményben $C\in\mathbb R$ tetszőleges integrációs állandó.

\begin{enumerate}[a)]
  \item Válasszuk $u=x$ és $\dd v=\cos x\dd x$ függvényeket. Ekkor
        $\dd u=\dd x$ és $v=\sin x$, tehát
        $$
          \boxed{\int x\cos x\dd x=x\sin x+\cos x+C}
          \text.
        $$

  \item Először legyen $u=x^2-1$ és $\dd v=\sin3x\dd x$. Ekkor
        $\dd u=2x\dd x$, $v=-\cos3x/3$, ezért
        $$
          \int(x^2-1)\sin3x\dd x
          =-\frac{x^2-1}{3}\cos3x+\frac23\int x\cos3x\dd x
          \text.
        $$
        A megmaradt integrálra újra parciálisan integrálunk:
        $$
          \int x\cos3x\dd x
          =\frac{x}{3}\sin3x+\frac19\cos3x
          \text.
        $$
        Így
        $$
          \boxed{\int(x^2-1)\sin3x\dd x
          =-\frac{x^2-1}{3}\cos3x
           +\frac{2x}{9}\sin3x+\frac{2}{27}\cos3x+C}
          \text.
        $$

  \item Az $1\cdot\ln x$ szorzatra alkalmazzuk a módszert ($x>0$):
        $$
          \boxed{\int\ln x\dd x=x\ln x-x+C}
          \text.
        $$

  \item Legyen $u=\arctan x$ és $\dd v=x\dd x$. Ekkor
        $\dd u=\dd x/(1+x^2)$ és $v=x^2/2$. Továbbá
        $x^2/(1+x^2)=1-1/(1+x^2)$, ezért
        \begin{align*}
          \int x\arctan x\dd x
          &=\frac{x^2}{2}\arctan x
            -\frac12\int\frac{x^2}{1+x^2}\dd x\\
          &=\boxed{\frac{x^2+1}{2}\arctan x-\frac{x}{2}+C}
          \text.
        \end{align*}

  \item Jelölje $I=\int e^x\sin x\dd x$. Kétszeri parciális integrálás után
        $$
          I=e^x\sin x-e^x\cos x-I
          \text,
        $$
        tehát
        $$
          \boxed{\int e^x\sin x\dd x
          =\frac{e^x}{2}(\sin x-\cos x)+C}
          \text.
        $$

  \item A $\sin^2x=(1-\cos2x)/2$ azonosságból
        $$
          \boxed{\int\sin^2x\dd x=\frac{x}{2}-\frac{\sin2x}{4}+C}
          \text.
        $$

  \item Legyen $t=\arccos x$, így $x=\cos t$ és
        $\dd x=-\sin t\dd t$. Ezután
        $$
          \int e^{\arccos x}\dd x=-\int e^t\sin t\dd t
          =\frac{e^t}{2}(\cos t-\sin t)+C
          \text.
        $$
        Visszahelyettesítve, $-1<x<1$ esetén
        $$
          \boxed{\int e^{\arccos x}\dd x
          =\frac{e^{\arccos x}}{2}
          \left(x-\sqrt{1-x^2}\right)+C}
          \text.
        $$
\end{enumerate}

\subsection{Racionális törtfüggvények}

\begin{enumerate}[a)]
  \item A nevező $(x-3)(x-4)$. Polinomosztás és parciális törtekre bontás
        után
        $$
          \frac{x^3-9x^2+27x-26}{x^2-7x+12}
          =x-2-\frac1{x-3}+\frac2{x-4}
          \text.
        $$
        Tagonként integrálva
        $$
          \boxed{\frac{x^2}{2}-2x-\ln|x-3|+2\ln|x-4|+C}
          \text.
        $$

  \item A nevezőben az $x$ háromszoros, az $x-2$ kétszeres gyöktényező.
        A felbontás
        $$
          \frac{x^3-2x^2+4}{x^3(x-2)^2}
          =\frac1{4x}+\frac1{x^2}+\frac1{x^3}
           -\frac1{4(x-2)}+\frac1{2(x-2)^2}
          \text.
        $$
        Ezért
        $$
          \boxed{\frac14\ln|x|-\frac1x-\frac1{2x^2}
          -\frac14\ln|x-2|-\frac1{2(x-2)}+C}
          \text.
        $$

  \item A nevezőt négyzetté egészítjük:
        $$
          x^2+4x+8=(x+2)^2+4
          \text,
        $$
        és a számlálót a nevező deriváltjához igazítjuk:
        $3x-2=\tfrac32(2x+4)-8$. Így
        $$
          \boxed{\int\frac{3x-2}{x^2+4x+8}\dd x
          =\frac32\ln(x^2+4x+8)
          -4\arctan\!\left(\frac{x+2}{2}\right)+C}
          \text.
        $$
\end{enumerate}

\subsection{Helyettesítéses integrálok}

\begin{enumerate}[a)]
  \item Legyen $t=\tan(x/2)$. Ekkor
        $$
          \cos x=\frac{1-t^2}{1+t^2},
          \qquad
          \dd x=\frac{2\dd t}{1+t^2}
          \text.
        $$
        Ezért
        \begin{align*}
          \int\frac{\dd x}{5+3\cos x}
          &=\int\frac{\dd t}{t^2+4}\\
          &=\boxed{\frac12\arctan\!\left(
            \frac12\tan\frac{x}{2}\right)+C}
          \text.
        \end{align*}
        Az eredmény az eredeti integrandus folytonossági intervallumain,
        megfelelő integrációs állandóval értendő.

  \item Legyen $u=\tanh(x/2)$. Ekkor
        $$
          \cosh x=\frac{1+u^2}{1-u^2},
          \quad
          \sinh x=\frac{2u}{1-u^2},
          \quad
          \dd x=\frac{2\dd u}{1-u^2}
          \text.
        $$
        Így
        $$
          \int\frac{\dd x}{1+\cosh x+2\sinh x}
          =\int\frac{\dd u}{1+2u}
          =\boxed{\frac12\ln\left|1+2\tanh\frac{x}{2}\right|+C}
          \text.
        $$

  \item A valós integrandus természetes tartományán $0<x<1$. Legyen
        $x=\sin^2t$, ahol $t\in(0,\pi/2)$. Ekkor
        $$
          \sqrt{\frac{x}{1-x}}=\tan t,
          \qquad
          \dd x=2\sin t\cos t\dd t
          \text.
        $$
        Ezért
        $$
          \boxed{\int\sqrt{\frac{x}{1-x}}\dd x
          =\arcsin\sqrt x-\sqrt{x(1-x)}+C}
          \text.
        $$

  \item A gyökkitevők nevezőinek legkisebb közös többszöröse $6$, ezért
        legyen $x=t^6$ ($t>0$). Ekkor
        $$
          \int\frac{\dd x}{\sqrt x(1+\sqrt[3]x)}
          =6\int\frac{t^2}{1+t^2}\dd t
          =6\int\left(1-\frac1{1+t^2}\right)\dd t
          \text.
        $$
        Visszahelyettesítve
        $$
          \boxed{6\sqrt[6]x-6\arctan\sqrt[6]x+C}
          \text.
        $$

  \item Legyen $t=e^x$, így $\dd x=\dd t/t$. Ekkor
        \begin{align*}
          \int\frac{e^x+2}{e^x+e^{2x}}\dd x
          &=\int\frac{t+2}{t^2(1+t)}\dd t\\
          &=\int\left(\frac1{t+1}-\frac1t+\frac2{t^2}\right)\dd t
          \text.
        \end{align*}
        Tehát
        $$
          \boxed{\ln(1+e^x)-x-2e^{-x}+C}
          \text.
        $$
\end{enumerate}

\end{document}
